%% =====================================================================
%%  附录 C　术语表
%%  编号体系独立于正文（附录 C / C.1 …）
%% =====================================================================
\chapter{术语表：人类智慧（反语）→ 大白话}

\applicable{你在某处看到一个词，不确定它到底在说什么。}

\section{使用说明}

值得注意的是，人类社会近年来出现了一批新的词。它们听起来是术语，实际上是在描述一些本来很朴素的事。

本表收录其中最常用的一批。\textbf{左列是它进场时的样子，右列是它出场时的样子。}

\begin{manstable}{P{3.2cm} L}
\tbh{术语} & \tbh{实际说的是} \\
\midrule
\textbf{边界感} & “这件事我不做，也不解释为什么。” \\
\textbf{情绪价值} & “我想让你听完之后感觉好一点，而不是知道这事该怎么办。” \\
\textbf{课题分离} & “这归你管，不归我管。” \\
\textbf{内耗／精神内耗} & “你在自己脑子里反复演，对面根本不知道。” \\
\textbf{松弛感} & “别紧张。” \\
\textbf{钝感力} & “有些话不用接。” \\
\textbf{配得感} & “这个东西我也可以要。” \\
\textbf{共情} & “我听见的不只是你说的那句话。” \\
\textbf{原生家庭} & “有些反应不是我选的。” \\
\textbf{允许一切发生} & “拦不住的就不拦了。” \\
\textbf{情绪稳定} & “他没在我面前崩。” \\
\textbf{断亲} & “我不打算再解释了。” \\
\textbf{搭子} & “只在一件事上跟你来往。” \\
\textbf{破防} & “那句正好戳到我了。” \\
\textbf{高能量} & “他愿意先开口。” \\
\textbf{显眼包} & “我不怕难看。” \\
\textbf{PUA} & “他在让我怀疑我自己。” \\
\textbf{有效沟通} & “我们互相都听懂了一点。” \\
\bottomrule
\end{manstable}

\section{关于这张表}

本表收录的词，有一个共同点：

\begin{mdquote}
\textbf{它们都在描述一件本来不需要名字的事。}
\end{mdquote}

人给一件事起名字，通常是因为这件事变难了。

饿了不需要“进食窗口管理”，困了不需要“睡眠优化”，走路不需要“位移方案”。\textbf{当“我不做这件事”需要被命名成“边界感”来辅助表达，说明直接说“我不做”，已经变得困难了。}

换句话说——

\begin{mdquote}
\textbf{这张表的长度，本身就是一份诊断书。}
\end{mdquote}

表越长，说明需要被翻译的东西越多。

而这本手册整本书，做的就是这样一件事：\textbf{把一种已经需要术语辅助才能成立的说法，重新翻译回它原来的样子。}

翻译本身的出现，说明原件已经不好认了。

%% 附录末「页脚交互条」由 hci-manual.cls 自动挂接，无需手写。
